\chapter{Research and Existing Solutions}
\section{Research Method / Approach}
We reviewed (i) peer-reviewed work on physical chess state reconstruction, (ii) open-source live-board projects, and
(iii) consumer products. For each we recorded the input (single image vs.\ stream), whether temporal or rule
information is used, the reported accuracy/latency, and its evaluation conditions. The review is ongoing; the novelty
statement in this chapter will be re-validated before FYP-2.
\todo{Add search log: databases/queries used, date, inclusion criteria.}

\section{Existing Solutions / Products}
\begin{description}
\item[Wölflein and Arandjelović] Board-state reconstruction from a single image with a RANSAC-based board
localiser and two CNNs, plus a large public dataset (roughly 1{,}450 participants and 2.5M frames in our notes). The
paper reports high per-square accuracy and low error on an unseen chess set \cite{wolflein2021}.
\todo{Verify exact figures and cite the dataset licence from the paper before submission.}
\item[Czyżewski et al.] Classical CV plus ML that iteratively processes the image; reported as lighting-resilient with high lattice-point and piece detection rates \cite{czyzewski2020}. \todo{verify numbers.}
\item[ARChessAnalyzer] A physical-board-to-analysis pipeline with CNN piece segmentation and augmented-reality output; our notes record about 93\% accuracy and a latency of several seconds from live capture to AR \cite{archess}. \todo{verify.}
\item[Wolarz] Automated chess move recognition \cite{wolarz}. \todo{summarise method and results.}
\item[Nayan] A tile-based recogniser operating on a live feed \cite{nayan}. \todo{summarise.}
\item[Consumer tools] Manual calculators (e.g.\ Chess Sigma Next Move) require the user to enter the position; post-game recorders (e.g.\ Chessgaze) digitise after play; smart boards provide live capture but cost $\$500+$.
\end{description}

\section{Comparison of Existing Solutions}
\begin{longtable}{p{0.2\textwidth}p{0.13\textwidth}p{0.13\textwidth}p{0.14\textwidth}p{0.13\textwidth}p{0.15\textwidth}}
\toprule
Solution & Live stream & Temporal reasoning & Legal-move constraints & Occlusion study & Needs special hardware\\
\midrule
Wölflein \& Arandjelović & No (image) & No & No\todo{verify} & Limited & No\\
Czyżewski et al. & No (image) & No & No\todo{verify} & Lighting only & No\\
ARChessAnalyzer & Yes (slow) & \todo{verify} & \todo{verify} & \todo{verify} & No\\
Nayan & Yes & \todo{verify} & Yes (legal inference) & \todo{verify} & No\\
Calculators / recorders & No & No & Partly & n/a & No\\
Smart boards & Yes & n/a & n/a & n/a & Yes ($\$500+$)\\
Chesstor (this work) & Yes & Yes & Yes & Evaluated (planned) & No\\
\bottomrule
\end{longtable}
Cells marked TODO are claims we have not yet verified against the primary sources; they must be completed
(the committee requires every claim to be specified and documented).

\section{Limitations and Gaps}
Prior systems are mostly evaluated as \emph{recognition} pipelines on per-image accuracy. What is under-studied, in
the sources we reviewed, is the \emph{state-integrity} problem in a live game: how often does a wrong or transient
observation corrupt the tracked game, how well can temporal and rule information prevent that, how quickly does
the system recover, and what does confirmation cost in latency. Existing live demos also tend to assume controlled
conditions, whereas the committee asks for robustness across angles, lighting, boards and piece materials.

\section{Proposed Solution}
We propose a streaming pipeline in which the trusted state changes only through a confirmation policy. Evidence is
accumulated in a temporal buffer; candidate transitions are \emph{legal move sequences} from the previous trusted
state; each is scored by agreement with the observed squares; occluded squares contribute no evidence; and a change is
committed only when the best hypothesis beats all others by a margin, otherwise the system waits or asks the user.
What we plan to do: build and evaluate this method against baselines on real boards. What we do \emph{not} claim:
novel board detection, novel piece classification, or a novel chess engine.

\section{Complex Computing Problem}
The problem has conflicting requirements (accuracy vs.\ latency vs.\ mobile compute; confidence threshold vs.\
responsiveness), no standard procedure for deciding whether a noisy change is a real move, interdependent stages
(localisation $\to$ recognition $\to$ inference $\to$ future state), ill-defined requirements (optimal window and
thresholds are empirical), and significant consequences (a wrong state invalidates all later analysis). The stress
experiments in Chapter~7 already show a genuine trade-off between recovering from persistent errors and
admitting false confirmations.

\section{Technical Challenges}
Hand and arm occlusion; motion blur and camera shake; glare and shadows; variation across boards and piece
materials (wood, glass, plastic); oblique camera angles; distinguishing similar pieces; special moves that change
more than two squares; correlated (persistent) recognition errors; real-time latency on a phone; and
reproducible evaluation without leaking board or session identity between training and test.

\section{Research Findings and Design Decisions}
\begin{itemize}
\item \textbf{Occupancy first:} the proposal POC asks whether occupancy/colour alone can identify the move. We built the core on occupancy+colour observations, with the rules supplying piece identity; full piece classification is added later and is needed for under-promotion.
\item \textbf{Legal hypotheses over diffing:} comparing against all legal sequences handles castling and en passant without special cases, where a naive source/destination diff cannot (Section~\ref{sec:results}).
\item \textbf{Occlusion as missing evidence:} occluded squares become ``unknown'', not ``empty'' (committee requirement).
\item \textbf{Temporal evidence by mean log-likelihood} over a window, which keeps decision margins on a single-frame scale and is robust to isolated glitches.
\item \textbf{Python core, then port:} the algorithm is developed and tested in Python first; the deployment path (on-device vs.\ local server) is decided only after profiling (Chapter~10).
\end{itemize}
